\section{Mathematical Model}

MANAR incorporates geographic navigation geometry, flight dynamics rules, and battery simulation models to govern aircraft behavior.

\subsection{Geographic Distance (Haversine Formulation)}
To evaluate spatial relationships over Earth's spherical surface without relying on planar approximations, MANAR calculates geographic distance using the Haversine formula.

Let two points on the Earth's surface be defined by their geodetic coordinates:
\begin{equation}
P_1 = (\phi_1, \lambda_1), \qquad P_2 = (\phi_2, \lambda_2)
\end{equation}
where $\phi$ represents latitude and $\lambda$ represents longitude in radians.

The angular differences are:
\begin{equation}
\Delta\phi = \phi_2 - \phi_1, \qquad \Delta\lambda = \lambda_2 - \lambda_1
\end{equation}

The square of half the chord length between the points is given by:
\begin{equation}
a = \sin^2\left(\frac{\Delta\phi}{2}\right) + \cos(\phi_1)\cos(\phi_2)\sin^2\left(\frac{\Delta\lambda}{2}\right)
\end{equation}

The angular distance in radians is:
\begin{equation}
c = 2 \operatorname{atan2}\left(\sqrt{a}, \sqrt{1 - a}\right)
\end{equation}

The surface distance $d(P_1, P_2)$ in meters is:
\begin{equation}
d(P_1, P_2) = R_E \cdot c
\end{equation}
where $R_E = 6\,371\,000.0\text{ m}$ is mean Earth radius. In \texttt{route\_optimizer.cpp} and \texttt{drone.cpp}, distance calculations use this exact double-precision Haversine formulation.

\subsection{Motion Model and Estimated Time of Arrival (ETA)}
The aircraft's horizontal velocity $v(m)$ is determined by the active operational preset $m \in \{\text{Quick}, \text{Active}, \text{Inspect}, \text{Hover}\}$:
\begin{equation}
v(m) = \begin{cases}
v_Q & m = \text{Quick} \\
v_A & m = \text{Active} \\
v_I & m = \text{Inspect} \\
v_H & m = \text{Hover}
\end{cases} \quad (\text{m/s})
\end{equation}
where $v_H = 0$ is defined semantically as stationary hover.

Given current aircraft coordinates $P_{\mathrm{drone}}$ and target coordinates $P_{\mathrm{dest}}$, remaining distance is $d = d(P_{\mathrm{drone}}, P_{\mathrm{dest}})$. For $v(m) > 0$, Estimated Time of Arrival ($t_{\mathrm{ETA}}$) is formulated as:
\begin{equation}
t_{\mathrm{ETA}} = \frac{d(P_{\mathrm{drone}}, P_{\mathrm{dest}})}{v(m)} \quad (\text{seconds})
\end{equation}

\subsection{Discrete Battery Simulation Model}
To exercise battery monitoring and failsafe behavior in the software prototype without requiring physical voltage/current telemetry, battery percentage $B_k \in [0, 100]$ evolves according to a discrete simulation rule:
\begin{equation}
B_{k+1} = \max\left(0.0, \; B_k - \alpha_m \cdot \Delta B\right)
\end{equation}
where $\Delta B > 0$ is a step decrement constant and $\alpha_m$ is a mode-dependent prototype simulation coefficient. When the aircraft is landed ($v = 0$, $H = 0$), battery consumption is paused ($\Delta B = 0$). This discrete simulation rule exists solely to exercise software battery and failsafe behavior, and is not a physical electro-aerodynamic battery model.

\subsection{Lawnmower Search Grid Geometry}
For a sector centered at $(\phi_c, \lambda_c)$ with sector width $W$ along longitude and height $H$ along latitude, conversion between meters and geographic coordinates uses local scale factors:
\begin{align}
M_{\mathrm{lat}} &= 111\,320.0 \quad (\text{m / deg lat}) \\
M_{\mathrm{lon}}(\phi_c) &= 111\,320.0 \cdot \cos\left(\frac{\phi_c \cdot \pi}{180.0}\right) \quad (\text{m / deg lon})
\end{align}

The southwest origin corner $(\phi_0, \lambda_0)$ of the grid is computed as:
\begin{equation}
\phi_0 = \phi_c - \frac{H / 2.0}{M_{\mathrm{lat}}}, \qquad \lambda_0 = \lambda_c - \frac{W / 2.0}{M_{\mathrm{lon}}(\phi_c)}
\end{equation}

Parallel sweep rows are separated by row spacing $\Delta y$.
